\section{Versionierte Migrationsfallstudie}
\label{chap:versioned-migration-case}

Der Migrationsfall beginnt mit einem älteren Toolingstand und seinem persistierten State.
Danach werden ausschließlich die austauschbaren Verzeichnisse \texttt{tools/} und
\texttt{docs/toolingdocs/} ersetzt. Eine registrierte Migration stimmt Konfiguration und
State auf die neue Version ab, ohne Produktdateien als Toolingbestand zu behandeln.

Die Prüfung unterscheidet gleichen Versionsstand, echten Versionssprung, Drift sowie
Fehler- und Rückgabefall. Die vorhandene historische Migrationssuite ist als \Verified{}
(\evidence{EV-UPGRADE-001}) referenziert; sie ist keine Behauptung, dass jede beliebige
Altversion unterstützt wird.
Abbildung~\ref{fig:tooling-migration} zeigt die Trennung von austauschbarem Payload
und erhaltenem Projektstate während der registrierten Migration.

\CaseDiagram{tooling-migration}{Migration vom alten Payload und State zu einem registrierten neuen Stand.}{fig:tooling-migration}
